% page 369 of the facsimile (image p-368 of the scan)
% block 165.1 x 237.7 mm ; left margin 36.9 mm
\pgc{15.240mm}{0.000mm}{
\l{\cel{54}361}
\l{}
\l{matrico. Mulduro qua produktas reliefajo per perkuto, giso,}
\l{\cel{3}e c. - DEFIRS.}
\l{}
\l{matrikulo. Listo, etato, en qua esas enskribita la nomi di omna}
\l{\cel{3}individui di korpo, di societo, - e, specale, listo, etato,}
\l{\cel{3}en qua onu enskribas la nomo, prenomo ed eniro-numero di}
\l{\cel{3}omna soldati qui eniras regimento.- DEFIRS.}
\l{}
\l{\sou{matrono.}\cel{2}Muliero respektinda pro lua matureso evala e lua bona}
\l{\cel{3}reputeso. - DEFIRS.}
\l{}
\l{\sou{matura}. I. (\sou{aludante frukto, grano}). Qua esas komplete desvol-}
\l{\cel{3}vita, developita.- II. (\sou{metaf}.)Qua jus atingis la punto maxim}
\l{\cel{3}favoroza (por), la maxim oportuna (por). - EFIS.}
\l{}
\l{\sou{matutino}. (\sou{liturgio katol.}) Parto unesma di la oficio, qua dic-}
\l{\cel{3}esis en la kloko unesma pos nokto-mezo - e dicesas nun (ecep-}
\l{\cel{3}te en ula monakeyi) en la vigilio. - EFIS.}
\l{}
\l{\sou{mauzoleo.}\cel{2}Tombo monumenta. - DEFIRS.}
\l{}
\l{\sou{maxilo.}\cel{2}(\sou{anat.}) I. La parto supra di la du parti osta di la}
\l{\cel{3}boko, qui suportas la denti che la animali vertebroza. -}
\l{\cel{3}II. (\sou{metaf.)}\cel{2}Nomo di la du fero-peci qui uzesas por tenar}
\l{\cel{3}fixe objekto per klemo. - EFI.}
\l{}
\l{\sou{maxim}. (\sou{gram.)}\cel{2}Indikas la gradaro-somito pri la eso-maniero}
\l{\cel{3}quan expresas adjektivo od adverbo. - DEFIRS.}
\l{}
\l{\sou{mayo}. La monato kinesma di la yaro. - DEFIRS.}
\l{}
\l{\sou{mayoliko}. Olima sorto di fayenco Hispana ed Italiana. - DEFIRS.}
\l{}
\l{\sou{mayonezo}. (\sou{koquarto}). Sauco kolda, ek oleo quan onu quirlas kun}
\l{\cel{3}vitelo til kande lu divenis mi-ferma. - DEFIRS.}
\l{}
\l{\sou{mayoro}. (\sou{armeo)}. Oficiro generala qua, helpe di generalisimo,}
\l{\cel{3}funcionas kom chefo di la stabo. - DEFIRS.}
\l{}
\l{\sou{mayuskulo}. Litero plu granda, e di formo maxim-multa-kaze}
\l{\cel{3}altra, kam la minuskulo korespondanta - uzata precipue kom}
\l{\cel{3}la rang-unesma litero di la nomi propra, e pos punto.DEFIS.}
\l{}
\l{\sou{mazo}. (\sou{olim}). Armo ek fero, kun kapo tre pezoza - en ula kazi:}
\l{\cel{3}garnisita ye pinti, quan onu manuagis quale Herkules manuagis}
\l{\cel{3}sua asom-bastono tuberoza, plu grosa ye un ek sua extremaji. -}
\l{\cel{3}EFIS.}
\l{}
\l{\sou{mazurko}.\cel{2}Danso segun mezuro tri-tempa, plu lenta kam valso,}
\l{\cel{3}ed en qua la tempo triesma esas plufortigita. - DEFIRS.}
\l{}
\l{me. (\sou{gram.)}\cel{2}Pronomo personala, qua reprezentas la persono qua}
\l{\cel{3}parolas o skribas. - DEFIRS.}
}
